\documentclass[11pt]{beamer}

% ===================== THEME & COLORS =====================
\usetheme{Madrid}
\usecolortheme{default}
\setbeamertemplate{navigation symbols}{}
\setbeamertemplate{footline}[frame number]

\definecolor{navy}{RGB}{27,42,74}
\definecolor{accentblue}{RGB}{47,84,150}
\definecolor{lightbg}{RGB}{235,240,248}

\setbeamercolor{structure}{fg=navy}
\setbeamercolor{title}{fg=white,bg=navy}
\setbeamercolor{frametitle}{fg=white,bg=navy}
\setbeamercolor{block title}{fg=white,bg=accentblue}
\setbeamercolor{block body}{fg=black,bg=lightbg}
\setbeamerfont{frametitle}{series=\bfseries}

\usepackage[utf8]{inputenc}
\usepackage[T1]{fontenc}
\usepackage{tcolorbox}
\tcbuselibrary{skins,breakable}
\usepackage{tikz}
\usetikzlibrary{positioning}
\usepackage{amsmath}
\usepackage{enumitem}

\newtcolorbox{infobox}[1]{
  colback=lightbg, colframe=accentblue, colbacktitle=navy,
  coltitle=white, fonttitle=\bfseries, title=#1,
  boxrule=0.4pt, arc=2pt, breakable
}

\newcommand{\circnum}[1]{%
  \tikz[baseline=-0.6ex]\node[circle,fill=navy,text=white,font=\bfseries\small,inner sep=1.6pt]{#1};%
}

% ===================== TITLE INFO =====================
\title[Maquette \(\rightarrow\) Code]{De la maquette au code\\[2pt]\large Génération automatique de code frontend multi-langage à partir d'une maquette UX haute fidélité}
\author[À COMPLÉTER]{KOUONGA TADIE KENNEDY}
\institute[Université de Douala]{
  Université de Douala\\
  Master II Recherche\\[4pt]
  \small Encadreur(s) : [À COMPLÉTER — Nom de l'encadreur]
}
\date{17 juillet 2026}

\begin{document}

% ------------------------------------------------------------
\begin{frame}
  \titlepage
\end{frame}

% ------------------------------------------------------------
\begin{frame}{Plan}
  \tableofcontents
\end{frame}

% ============================================================
\section{Problématique}
\begin{frame}{Problématique}

  \vspace{0.5em}
  \begin{infobox}{Question centrale}
    Comment automatiser la transformation d'une maquette UX haute fidélité en code
    frontend fonctionnel, fidèle visuellement, et généralisable à plusieurs langages,
    afin de recentrer le développeur sur la logique métier ?
  \end{infobox}
\end{frame}

% ============================================================
\section{Objectif global}
\begin{frame}{Objectif global}
  \begin{infobox}{Objectif}
    Concevoir une méthode capable de transformer automatiquement une \textbf{maquette
    UX haute fidélité} en \textbf{code frontend} fonctionnel, fidèle au design
    d'origine, et \textbf{généralisable à plusieurs langages / frameworks} cibles.
  \end{infobox}
  \vspace{1em}
  \begin{center}
  
  \end{center}
  \begin{center}
  \end{center}
\end{frame}

% ============================================================
\section{Objectifs spécifiques}
\begin{frame}{Objectifs spécifiques}
  \begin{itemize}
    \setlength{\itemsep}{0.6em}
    \item[\textbf{1.}] \textbf{Identifier} les  métriques mesurables d'une
      maquette haute fidélité (grille, marges, palette, échelle typographique,
      positions des éléments).
    \item[\textbf{2.}] \textbf{Concevoir} un modèle structurel intermédiaire (qui est la representation hierachique sous forme d'arbre des differents elements de notre maquette) représentant fidèlement l'organisation de la maquette.
    \item[\textbf{3.}] \textbf{Déterminer} une méthode permettant d'établir les liaisons et les relations entre éléments dans notre model de façon fiable.
    \item[\textbf{4.}] \textbf{Développer} un mécanisme de génération de code à partir de ce
      modèle, adaptable à plusieurs langages de programmation ou framework.
    \item[\textbf{5.}] \textbf{Évaluer} la fidélité et la qualité du code généré par rapport à la
      maquette d'origine.
  \end{itemize}
\end{frame}

% ============================================================
\section{Questions de recherche}
\begin{frame}{Questions de recherche}
  \begin{itemize}[leftmargin=1.4em, itemsep=0.9em]
    \item[\circnum{1}] \textbf{QR1} — Quelles  métriques
      permettent de décrire objectivement une maquette haute fidélité ?
    \item[\circnum{2}] \textbf{QR2} — Comment transformer une maquette (données
      visuelles) en un modèle structuré  ?
    \item[\circnum{3}] \textbf{QR3} — Quelles sont les differentes methodes permettent de déterminer fiablement les liaisons entre les
      éléments d'une maquette ?
    \item[\circnum{4}] \textbf{QR4} — Comment, à partir de ce modèle, générer du
      code frontend fonctionnel et fidèle, dans différents langages ?
  \end{itemize}
  \vspace{0.8em}
\end{frame}

% ============================================================
\section{Développements réalisés}
\begin{frame}{Résumé — de la maquette aux metriques mesurables}
  \small
  \begin{block}{Caractérisation de la maquette}
    \begin{itemize}
      \item Dimension du canevas, marges, rythme vertical (grille de 8 px)
      \item Palette de couleurs limitée et réutilisée
      \item Échelle typographique restreinte (titre, sous-titre, corps)
     
    \end{itemize}
  \end{block}
  \begin{infobox}{Pourquoi c'est la première étape}
    Ces métriques transforment une image (pixels) en un ensemble d'\textbf{éléments
    mesurables} (position, taille, style) qui serviront de base à la construction d'un modèle structurel.
  \end{infobox}
\end{frame}

\begin{frame}{determination des liaison entre les differents elements(methode de similarite)}
  \small
  Chaque paire d'éléments est comparée par un score combinant proximité spatiale et
  ressemblance visuelle en ce basant sur les metriques qu'on a eu a detecter dans la maquette. La formule de similarité est donnée par :
  \[
    \sigma(b_i,b_j) = \alpha \cdot \text{sim}_{\text{geo}}(b_i,b_j) + \beta \cdot \text{sim}_{\text{sem}}(b_i,b_j)
  \]
 
\end{frame}

% ============================================================\section{Conclusion et perspectives}
\begin{frame}{Conclusion}
  \begin{infobox}{Ce qui a été établi}
    Nous avons défini les différentes métriques d'une maquette
    haute fidélité, et, dans le but de concevoir notre modèle, nous avons étudié la
    relation entre les éléments de notre maquette UX à travers une méthode de calcul
    de similarité.
  \end{infobox}
  \vspace{0.6em}
  Cette étape pose les fondations nécessaires à la construction du modèle structurel
  puis à la génération automatique de code frontend.
\end{frame}

\begin{frame}{Perspectives}
  \begin{itemize}[leftmargin=1.2em, itemsep=0.7em]
    \item \textbf{Explorer d'autres méthodes} de calcul de similarité proposées
      dans la littérature (différents articles de recherche).
    \item \textbf{Comparer ces méthodes} entre elles et adapter, dans un premier
      temps, la plus pertinente à notre cas d'étude.
    \item \textbf{Déterminer notre modèle} définitif à partir de la méthode retenue.
    \item \textbf{Étudier la génération de code} : comment produire, à partir de ce
      modèle, plusieurs versions de code dans différents langages.
  \end{itemize}
\end{frame}

\begin{frame}
  \begin{center}
    {\Huge\bfseries\color{navy} Merci}\\[1em]
    {\large Questions ?}
  \end{center}
\end{frame}

\end{document}